\documentclass[10pt,letterpaper,twocolumn]{article}
\usepackage[letterpaper,top=0.65in,bottom=0.65in,left=0.66in,right=0.66in,columnsep=0.22in]{geometry}
\usepackage{times}
\usepackage[T1]{fontenc}
\usepackage{microtype}
\usepackage{enumitem}
\usepackage{url}
\usepackage[hidelinks]{hyperref}
\usepackage{titlesec}
\setlength{\parindent}{1em}
\setlength{\parskip}{0pt}
\setlist{nosep,leftmargin=*}
\titleformat{\section}{\bfseries\normalsize}{\thesection.}{0.45em}{}
\titleformat{\subsection}{\bfseries\small}{\thesubsection.}{0.4em}{}
\titlespacing*{\section}{0pt}{0.7ex}{0.3ex}
\titlespacing*{\subsection}{0pt}{0.55ex}{0.2ex}
\pagestyle{plain}

\begin{document}
\twocolumn[
\begin{center}
{\Large\bfseries Artifact Appendix: Versatile yet Efficient Network Traffic Analysis: Offloading Network Foundation Model to SmartNIC}\par
\vspace{0.45em}
{\small Badges requested: Available and Functional (Reproduced is not requested)}
\end{center}
\vspace{0.5em}
]

\section{Artifact Overview}
This artifact provides the implementation and evaluation code for \textsc{Nepco}, the network foundation model presented in the paper. The release retains the complete workflows for PCAP processing, pre-training and fine-tuning, host and SmartNIC inference, ONNX conversion, and adversarial evaluation. It also contains two bounded smoke tests that require neither retraining nor external traffic captures: E1 reproduces DataCon classification metrics from a packaged fine-tuned checkpoint and labeled TSV, and E2 verifies the CPU inference-timing path using the same artifacts.

\section{Description and Requirements}
\subsection{Artifact Layout}
The implementation, configuration, vocabulary, and pre-trained model are under \texttt{uer/}, \texttt{configs/}, \texttt{vocab/}, and \texttt{models/}. Complete workflows are under \texttt{scripts/}. The two self-contained checks are under \texttt{E1/} and \texttt{E2/}; each has a README and manifest. E1 additionally contains the fine-tuned checkpoint and test TSV used by both checks.

\subsection{Hardware Dependencies}
E1 requires only an x86-64 Linux host and does not require a GPU. E2 additionally requires Linux \texttt{taskset} and exactly eight available logical CPU cores. We recommend 8~GB RAM and 1~GB free disk space for these smoke tests. Full model training and SmartNIC deployment have additional GPU and BlueField-3 requirements documented in the top-level README and are outside the bounded E1/E2 evaluation.

\subsection{Software Dependencies}
Python dependencies are listed in \texttt{requirements.txt}. The verified environment uses Python~3.10.14 and PyTorch~2.3.0. E2 uses the standard Linux CPU-affinity interface through \texttt{taskset}. No network access is required after the artifact and Python dependencies have been installed.

\subsection{Packaged Inputs}
E1 packages \path{E1/artifacts/DataCon/finetuned_model.bin} and \path{E1/artifacts/DataCon/test_dataset.tsv}. The TSV contains 2,857 labeled test samples plus its header. Their sizes and SHA-256 hashes are recorded in \path{E1/manifest.json}. E2 references these files instead of shipping a duplicate checkpoint.

\section{Installation}
From the repository root, create an isolated environment and install the dependencies:
\begin{verbatim}
python3 -m venv .venv
source .venv/bin/activate
python3 -m pip install \
 -r requirements.txt
\end{verbatim}
Confirm that both E1 input artifacts, \texttt{configs/nepco\_config.json}, and \texttt{vocab/hex\_vocab.txt} are present. Both evaluators load the fine-tuned state dictionary with \texttt{strict=True}; an architecture mismatch therefore fails rather than silently testing a partially initialized model.

\section{Functional Claims}
\begin{itemize}
\item \textbf{C1---Classification path.} The packaged final Nepco architecture, fine-tuned DataCon checkpoint, vocabulary, and labeled TSV can be loaded together to reproduce the archived classification output. E1 tests this claim on all 2,857 samples.
\item \textbf{C2---CPU timing path.} The same packaged model and input representation can execute through the CPU latency path with batch size~1 on eight fixed logical cores. E2 tests this claim with a bounded eight-sample run.
\end{itemize}
These checks validate executable release paths. They do not retrain Nepco, reproduce every paper table, compare against ET-BERT, or exercise the BlueField-3 testbed.

\section{Evaluation}
\subsection{E1: DataCon Classification Smoke Test}
\textbf{Effort:} less than one human-minute and typically less than one compute-minute on a recent CPU. Run:
\begin{verbatim}
PYTHON_BIN=python3 \
 bash E1/run_smoke_test.sh
\end{verbatim}
The evaluator loads \texttt{finetuned\_model.bin}, reads the complete packaged TSV, and writes \texttt{metrics.json}, per-class \texttt{prf.csv}, and \texttt{predictions.tsv} under \path{E1/results/DataCon/}. The verified output is macro precision~0.992038, macro recall~0.991294, macro F1~0.991604, and accuracy~0.986699 over 2,857 samples. Successful strict checkpoint loading, completion of all samples, and agreement with the archived files validate C1.

\subsection{E2: Nepco CPU Latency Smoke Test}
\textbf{Effort:} less than one human-minute and typically a few compute-seconds after model loading. Select any eight logical CPU cores available to the host and run:
\begin{verbatim}
PYTHON_BIN=python3 CPU_CORES=0-7 \
 bash E2/run_smoke_test.sh
\end{verbatim}
E2 reuses the E1 checkpoint and TSV, processes the first eight samples with batch size~1, performs one warm-up pass, and records five measurement passes. Only \texttt{model.infer} calls are timed; checkpoint loading, tokenization, tensor construction, and output serialization are excluded. The script writes \texttt{nepco.json} and \texttt{nepco.csv} under \path{E2/results/DataCon/}.

On the reference host, using cores 168--175, E2 completed at $2.271679 \pm 0.201898$~ms/sample (95\% confidence-interval half-width). This host-specific value confirms that the timing path works; it is neither a full-dataset result nor a cross-system comparison. Successful completion and correctly recorded affinity validate C2.

\section{Customization and Interpretation}
Reviewers may append \texttt{--max\_samples 32} to E1 for a shorter classification code-path check, although the archived metrics use all 2,857 samples. E2 accepts \texttt{MAX\_SAMPLES} and \texttt{CPU\_CORES}; the latter must select exactly eight logical cores. Changing the sample count, core placement, software stack, or host invalidates numerical latency comparison but not the functional criterion.

Full pre-training, fine-tuning, multi-dataset evaluation, ONNX conversion, BlueField-3 deployment, and adversarial robustness remain documented in the top-level README and \texttt{scripts/}. They are intentionally outside these bounded E1/E2 smoke tests.

\end{document}
